\chapter{Another Hello to the Third Man}
\chaptermark{Another Hello to the Third Man}
\chapterauthor{Julian Torres, \textit{
University College London
}}

\begin{quote}
    This essay presents some objections to a prominent solution to the Third-Man Argument (TMA) as it appears in the Parmenides due to Constance Meinwald (2022). According to this solution, the regress seemingly generated by the TMA can be blocked by distinguishing two kinds of predication: Ordinary Predication, which attributes properties to particulars, and Tree Predication, a special form of self-predication that specifies the nature or definition of a Form.

I argue that Meinwald’s Tree Predication model ultimately fails. First, I argue that the distinction between Ordinary and Tree Predication does not successfully generalise to cases discussed in the Sophist, such as predications involving the Form of the Same, without reintroducing a version of the regress it was meant to avoid. Second, I argue the Tree Predication model lacks adequate explanatory power: when multiple Tree Predications are available, the view yields explanatory overdetermination, weakening the original motivation for the account. Once I suggest some ways Meinwald may respond to these objections, I argue her resulting position either reintroduces a regress at the level of Forms or deprives the Tree-Predication model of its explanatory power and philosophical motivation. 
\end{quote}

\section{Introduction}
In the Sophist, the Visitor introduces a seemingly absurd thesis that `young people and old late-learners’ hold: \textit{only the good is good and the man is man} (Sph, 251c). Surprisingly, it seems uncontroversial to everyone involved in the discussion that predications of the form `The Man is man’ are intelligible. But understanding the way in which Plato thought of these instances of so-called Self-Predication has proven to be no easy task. It has also proven to be crucial to reconstructing his views on the Third Man Argument of the Parmenides. 

By an examination of evidence from the Sophist and the Parmenides, I evaluate one prominent view in the literature by Constance Meinwald of Plato's views on Self-Predication in the Parmenides. First, I reconstruct the Third-Man Argument (TMA) as found in the \textit{Parmenides} and assess the success of Meinwald's proposal in avoiding the TMA's conclusion. After this, I raise some objections to the proposed solution. I argue that, although it correctly identifies Self-Predication as the premise to qualify in order to avoid the infinite regress of Forms that gives the TMA its argumentative weight, the model of Tree Predication that Meinwald defends is ultimately unsuccessful as a solution to the TMA, and thus an inadequate reconstruction of Plato's own views. However, in order to get a clearer picture of the place of Meinwald's solution in the overall dialectic, some familiarity with Plato's theory and the challenge posed to it by the TMA is required. 

Plato's forms play crucial roles in his metaphysics and epistemology. In general, they appear to be abstract entities, as opposed to sensible ones, which possess at least some subset of the following properties: they are causally responsible for and explanatory of the properties of sensible objects, they are the objects of knowledge, and grasp of them seems to lead to ethical improvements in our lives in some substantive sense (Harte, 2019, 421--422). 

Now, the Third-Man argument, as traditionally construed, purports to undermine the coherence of this broad notion of the Theory of Forms. Thus, if successful, the TMA seriously compromises the status of one of the central pieces of Plato's broader philosophy. I will now reconstruct the argument, and then consider one suggested solution to it. 

\section{The Third Man}

The Third-Man Argument is fairly short, and it goes as follows (\textit{Prm}, 132a--132b):

\begin{quote}
\textbf{P1. One-over-many}: Let $A_1, \dots , A_n$ be some objects. If $A_1, \dots, A_n$ are all $F$, then there is a thing, F-ness, such that F-ness is a unique Form in respect of which A$_1$, … A$_n$ are all $F$ (132a). \\
\textbf{P2. Self-Predication}: For all Forms F-ness (Largeness, Smallness, etc.), `F-ness is F' holds. (\textit{Sph}. 258b9-c3). \\
\textbf{P3. Non-Identity}: For all $x$, if $x$ is $F$, then $x$ is not identical with F-ness (This premise is supplied following Vlastos (1954, 325)).
\end{quote}

\noindent Consider now the form of Largeness and a set of large things. From \textbf{One-over-many} and \textbf{Self-Predication}, it now follows that Largeness is large and that there is a unique Form of \textit{Largeness} in respect of which Largeness and the members of the set are large. But from \textbf{Non-Identity} it follows that Largeness is not identical to the Form of \textit{Largeness}. So Largeness and other large things must be large with respect to a new form, Largeness$_1$. By iterating the argument \textit{ad infinitum}, it follows that:

\begin{quote}
\textbf{C.} There is an infinite set of Forms {Largeness, $Largeness_1$, … $Largeness_n$}, where $Largeness_k$ explains the property of being large for all $n<k$. 
\end{quote}

\noindent Informally, we have an infinite regress of Forms of Largeness. The regress seems vicious in at least one sense. If F-ness is meant to explain the features of things that are F, then one ought to first understand how F-ness$_1$ explains the features of F-ness, and so on. Thus, F-ness seems to never fulfil its intended explanatory role. 

One option to reject C is to note that it follows from a contradiction: Self-Predication and Non-Identity are jointly inconsistent (Vlastos 1954, 326) so one (or both) of them must be rejected or qualified to not contradict each other. In what follows, I present a modification of Self-Predication by Meinwald (2022) that distinguishes between Ordinary Predication and predication of the form `Largeness is large'.

\section{Self-Predication}

Consider the following instance of Self-Predication (\textit{Sph}, 251c): 
\begin{quote}
M: The Man is a man.
\end{quote}

\noindent On a natural reading of the sentence, `The Man’ refers to the Form in virtue of which all men are men. By \textbf{One-Over-Many}, if Plato and Aristotle are both men, then there is a Form `the Man’, and so on.\footnote{And we want to prevent the need for a  `Third-Man' to make sense of the predication!} But, if we follow the natural reading suggested, the predicate term which applies to the Form `The Man’ is left unclear. Quite plausibly, Plato is not committing himself to the view that the form of Man-ness walks around Athens, having conversations with other men. The Man must be a man in some way different from how Aristotle is a man. 

\subsection{Ordinary Predication}

It is worth clarifying Meinwald's position further. In her view, Self-Predication stands in contrast with \textbf{Ordinary Predication}. In cases such as `$x$ is $F$', where $F$ is some property and $x$ is some particular thing, $F$ simply stands for a particular feature or property which $x$ has. In the TMA, `$x$ is large' is \textbf{ordinary} when it just means what it says: that $x$ is a large thing.

\subsection{Tree Predication}

On one construal (Meinwald 2022 and 1992), Forms stand in predication with themselves in virtue of a relation that differs from the relation that obtains between a subject and its property. On this view, Plato uses the Third Man argument to motivate a distinction that will be important in the second half of the \textit{Parmenides}, namely that of predication `in relation to itself’ as distinct from predication `in relation to others’. The former is called \textbf{Tree Predication (TP)}, and the latter \textbf{Ordinary Predication (OP)}.  When in the predication `$x$ is $F$,’ $x$ is the property corresponding to F-ness, then $F$ in `is $F$' does not express a property of $x$ but a part of the definition of what it is to be $x$. Given that instances of predication in which concrete particulars figure seem to express properties instead of definitions predication `in relation to itself’ only occurs when $x$ is a term referring to a Form.\footnote{One can define concrete particulars in classical predicate logic in a suitable model, but these definitions do not take the form of self-predication, so my point stands.}

Applying the distinction to the Third Man Argument, Meinwald’s view amounts to a rejection of Non-Identity when $x$ is a Form and `is $F$' is predicated of $x$; no third form is needed because `Largeness is large' just specifies the nature of largeness. 

The view strikes me as plausible in one important way: it maintains Self-Predication while retaining the intuitive principle of Non-Identity for Ordinary Predication between a particular thing and its properties. Indeed, it is evident throughout the \textit{Sophist} that \textit{there is} in fact a difference between the features which things have `in virtue of their own nature' or `in virtue of sharing in some other (Form)' (255e; 250c). When `the apple is red' holds, the apple is certainly not identical with the \textit{Redness} in virtue of which it is red. But in `Redness is red', the subject and the form \textit{are identical}. This gives us reason to think that a different kind of relation between subject and predicate holds in each case. I shall assess the plausibility of Meinwald's distinction by applying it to the discussion of kinds in the \textit{Sophist} (250a--258d). If the distinction cannot account for the various ways in which certain Forms are said to stand in relation to themselves and other Forms, then we will have reason to doubt the success of the distinction in dealing with the TMA. 

\section{The \textit{Sophist}}

So far, the proposed solution to the TMA distinguishes Tree from Ordinary Predication, thereby weakening the Non-Identity premise of the argument and avoiding the infinite regress of forms. But now consider the following predication implicit in the \textit{Sophist} (256a)\footnote{To be sure, one may offer a solution to the TMA argument as it arises in the Parmenides only. Such a solution need not, presumably, account for instances of self-predication in other dialogues. Thankfully, Meinwald (2022, p. 427) makes clear that her model of Tree-Predication is intended as a model for instances of self-predication in the Sophist too. Thus, by her own lights, failure to avoid the TMA as it arises in the Sophist bears on the plausibility of her solution in the Parmenides too.}:
\begin{quote}
\textbf{S.} The Same is the same.
\end{quote}

\noindent Following Meinwald, we may want to distinguish in this case both senses of predication present in \textbf{S}. In this case (noting that `$x$ is the same' means $x$ is the same as itself, i.e., $x$ is self-identical), Plato seems to be committed to: 
\begin{quote}
\textbf{Ordinary Predication (OP$_S$)}: The Same is the same.
\end{quote}

\noindent Note that the weakened form of Non-Identity which emerged from Meinwald's solution of the TMA applied only to Ordinary Predication. If this is right, then one can immediately see how Meinwald's solution fails to avoid the infinite regress in cases like OP$_S$. One may run a version of the TMA where The Same and a collection of other things (anything whatsoever, because everything is self-identical) are said to be `the same', and an infinite regress of Forms The Same$_1$\dots The Same$_n$ will emerge to explain the relevant instances of Ordinary Predication. If instances of self-predication are meant to capture a definitional kind of predication for forms, Meinwald must explain cases in which forms, as entities, possess properties and may self-predicate in the \textbf{OP} sense.

However, Meinwald considers cases like \textbf{OP$_S$} to be unproblematic (Meinwald, 2022, 424--425). In her view, the solution to this new bolstered version of the TMA is simple. No new Form of The Same apart from the original one will emerge because there is no explanatory need for it. When various things, including the Form of The Same itself, are said to have the property of being the same (or self-identical), then this feature is explained by what she calls \textit{Platonic Anaxagoreanism} (Meinwald, 2022, 411). According to this view, particular things' properties are explained by the features of their metaphysical constituents or ingredients. In particular, if `$x$ is the same (as itself)' is true, then this is explained by $x$ having a share of The Same. But, so the story continues, it follows from Self-Predication in Meinwald's interpretation that:
\begin{quote}
\textbf{Tree Predication (TP$_S$)}: The Same is the same as itself.
\end{quote}

\noindent And it is precisely \textbf{TP$_S$} that explains why having a share of The Same makes something self-identical (Meinwald, 2022, 419--420). In other words, no extra form of The Same is needed because \textbf{TP$_S$} explains \textbf{OP$_S$}. The form of The Same has the property of being self-identical in virtue of itself. 

Note that this amounts to a rejection of Non-identity for certain cases of Ordinary Predication. However, Meinwald (Meinwald, 2022, 424--425) thinks she has a principled reason for doing so: Forms have explanatory roles which no sensible particular can perform; it is fine for Forms to explain their own properties and those of other particulars because they are the only things which could play such roles.

But it seems to me that this response is not satisfactory. One could equally run a bolstered version of the argument where every subject of predication is itself a Form (for Forms are also self-identical). Consider the set of \textbf{OP} predications: 
\begin{quote}
\textbf{Ordinary:} The Same is the same, The Fast is the same, The Slow is the same.
\end{quote}

\noindent In the case of \textbf{Ordinary}, the response from the unique explanatory and causal roles that forms perform is no longer available to Meinwald. Any of the Forms (The Fast, The Slow, or The Same) can play the relevant explanatory role if being a Form is the only requirement that matters. Indeed, if Meinwald is to assert that the Form of The Same has a distinct status \textit{qua} Form to explain instances of the \textbf{OP} predicate: `is the same', then more than merely meeting the requirement of being a form is needed. 

Indeed, Meinwald does seem to think that instances of tree-predication for the Form of The Same provide the explanatory ground for the ordinary predication of objects which exhibit sameness (Meinwald, 2022, 421), and thus, she thinks that there is reason to think that the Form of The Same will take precedence in providing the relevant ground for the property of sameness which particulars exhibit.\footnote{I return to this issue in Section 5}

\subsection{Explanatory Power}

In any case, there are further arguments available to cast suspicion over Tree Predication's explanatory power. Consider the following cases of \textbf{TP} and the correlative cases of \textbf{OP}, all of which Meinwald accepts (1992, p.383):
\begin{quote}
\textbf{TP$_B$}: The Bear is omnivorous.\\
\textbf{OP$_B$}:  Bear $x$ is omnivorous.\\
\textbf{TP$_X$}: The Fox is omnivorous.\\
\textbf{OP$_X$}: Fox $x$ is omnivorous.
\end{quote}

\noindent Meinwald's view seems to be that each of the \textbf{OP} cases is explained by the relevant \textbf{TP} (assuming, of course, that Bear $x$ has a share of bear-hood): The bears in the Yosemite Park eat plants and meat because part of what it is to be a bear is to eat those things, and the same goes for foxes. But it also seems to follow, by \textbf{One-Over-Many}, that there is a common property which both foxes and bears possess by having a share in some Form of Omnivorousness, and which in turn explains their feature. Now, Meinwald must explain which one of these Forms is meant to be the truly explanatory one of the properties the particulars have. She should explain which instance of Tree Predication grounds the sensible property: 
\begin{quote}
    \textbf{TP$_X$, TP$_B$}
\end{quote}

 \noindent or 
\begin{quote}\textbf{TP$_O:$} Omnivorousness is omnivorous. 
\end{quote}

\noindent Without a more complete story, Tree Predications seem to be too broad to give the precise explanatory ground required by Ordinary Predications. Instead, they result in explanatory overdetermination for sensible properties in cases where the grounding is clear---it seems plausible that Omnivorousness explains the diet of bears! More is needed to account for the explanatory priority of some forms over others when their Tree Predications overlap. Thus, I consider a way in which Meinwald may tweak her view to account for cases of overdetermination. The upshot of it is a much weaker case for accepting the Tree Predication solution.

\section{Another regress?}

Consider again the case above. In order to account for the priority of Omnivorousness in providing the explanatory ground for \textbf{OP$_X$} and \textbf{OP$_B$}, Meinwald may ask us to run the following test. Whenever more than one Tree Predication seems a plausible candidate to explain the same instance of Ordinary Predication, then one must choose the most specific one to the nature of the respective form. Consider: 
\begin{quote}
\textbf{TP$_B$}: The Bear is omnivorous.\\
\textbf{TP$_O$}: Omnivorousness is omnivorous.
\end{quote}

\noindent In this case, we must choose \textbf{TP$_O$}, for being omnivorous is the most specific \textbf{TP} possible for Omnivorousness, whereas it is not for The Bear. Being omnivorous fully specifies the nature of Omnivorousness, whereas many other \textbf{TP}s must be predicated of the Bear for a full definition. Bears share the feature of being omnivorous with other species distinct from them, and it seems plausible that a full definition of bear-hood ought to reflect that fact by providing distinctions that are fine-grained enough and do not assign the same definition to bears as it does to any other omnivorous animals. For instance, it seems that \textbf{TP}s specifying that The Bear is mammal, of the family \textit{Ursidae}, and so on will belong to a full definition of The Bear via \textbf{TP}s. However, a simpler view is possible. By iterating the definitional process above, it seems one may reach a maximally specific instance of \textbf{TP} for the form of The Bear which no other animal form may share, namely:
\begin{quote}
\textbf{TP$_{Bear}$}: The Bear is a bear.
\end{quote} 

\noindent Let instances of \textbf{TP} that are maximally specific like \textbf{TP$_{Bear}$} be known as \textbf{Maximal Tree Predication (MTP)}. The upshot of this modified version of the view with \textbf{MTP}s may go as follows. The required explanatory ground for any \textbf{OP} of the form `the $x$ is $F$' is provided by the \textbf{MTP} of F-ness (i.e., the $F$ is $F$). And so it seems we now obtain the required ground: Bear $x$ is omnivorous in virtue of Omnivorousness, whose nature is just to be omnivorous.

Now, it seems that the source of the explanatory overdetermination in cases like the above is that `The Bear is omnivorous' holds in virtue of a form distinct from it. Even though it is  in the definition of The Bear to be omnivorous, this fact does not provide the required explanatory ground. We require a second form. To bring out the consequences of this version of Meinwald's view, consider a bolstered version of the Third Man argument where every subject of predication is a form. Let the subscript `$\star$' stand for an instance of \textbf{TP}:
\begin{quote}
\textbf{P1. One-over-many$_*$}: Let $A_1,\dots  , A_n$ be some forms. If $A_1,\dots , A_n$ are$_*$ all $F$, then there is a form, F-ness, such that it is a unique Form in respect of which A$_1$, … A$_n$ are$_*$ all $F$\dots 

And so on. 
\end{quote}

\noindent In order to stop this regress, a similar response is available to Meinwald: no Third Form arises from the argument because F-ness is $_*$ $F$ for it is an instance of (\textbf{MTP}) in a different way to how X-ness is$_*$ $F$. 

Now, this response may be successful in stopping the regress, but it appears to undermine the original motivation to accept Meinwald's account of Self-Predication. On her view, instances of \textbf{TP} ground the sensible properties of particulars because they specify the metaphysical ingredients of the Forms. Yet, the needed modifications to her view via \textbf{Maximal Tree Predications} do away with the explanatory power of any instance of \textbf{TP} that is not \textbf{MTP}. If Platonic Anaxagoreanism is right, it no longer seems that the metaphysical constituents of forms are specified by the \textbf{TP}-model, and the view that Self-Predication is Tree Predication loses much of its plausibility. Unlike in Meinwald's original proposal, on the \textbf{MTP}-model, instances of self-predication do not capture part of the definition of what it is to be some form, and so most instances of \textbf{TP} seem both philosophically unmotivated and explanatorily inert. If this is right, most of Meinwald's \textbf{Tree-Predication} model can be done away with. 

\section{Conclusion}

In sum, I assessed a proposal to solve the TMA by Constance Meinwald in light of a passage of the \textit{Sophist} that features instances of Ordinary Predication for forms which may be problematic for her solution. I offered some objections to her view and suggested ways in which her view may be modified to avoid them. I have argued that the modifications required undermine the original motivation for her solution and make much of the structure it required for her solution unnecessary. 

\refsection

\begin{hangparas}{\hangingindent}{1}
Cooper, John, ed. 1997. \textit{Plato: Complete Works}. Indianapolis: Hackett.

Harte, Verity. 2019. ``Plato's Metaphysics.'' In \textit{The Oxford Handbook of Plato}, edited by Gail Fine, 2nd ed. Oxford: Oxford University Press.

Meinwald, Constance. 1992. ``Good-bye to the Third Man." In \textit{The Cambridge Companion to Plato}, edited by Richard Kraut. Cambridge: Cambridge University Press.

Meinwald, Constance. 2022. ``Another Goodbye to the Third Man.” In \textit{The Cambridge Companion to Plato}, edited by David Ebrey and Richard Kraut. Cambridge: Cambridge University Press. 

Vlastos, Gregory. 1954. ``The Third Man Argument in the \textit{Parmenides}.” \textit{The Philosophical Review} 63 (3): 319--49. 
\url{https://doi.org/10.2307/2182692.}
\end{hangparas}